\section{The Normal \texorpdfstring{\(\Theta\)}{Theta}: Origin, Geometry, and Computability}
\label{sec:normal-theta}

This section supplies the geometry of the motivating spectral-level-set
application.  It answers four distinct questions: what motivates a spectral
scale, which extra modeling choice creates an exact level set, why its normal is
\(\Theta=u_1v_1^\top\), and what can actually be certified when the leading
singular vectors are computed approximately.

\subsection{Spectral scaling is the motivation, not the equality constraint}

Let \(W\in\R^{m\times n}\) represent a hidden linear map \(y=Wx\).  Since
\[
 \frac{\|Wx\|_2/\sqrt m}{\|x\|_2/\sqrt n}
 \le \snorm{W}\sqrt{\frac nm},
\]
the spectral formulation of width-stable feature learning motivates the scale
\citep{yang2023spectral}
\begin{equation}
 \snorm{W},\ \snorm{\Delta W}
 =\Theta\!\left(\sqrt{\frac mn}\right)
\tag{N1}
\label{eq:spectral-scaling}
\end{equation}
Equation \eqref{eq:spectral-scaling} is an asymptotic order statement.  It does
\emph{not} require the exact equality \(\snorm{W}=R\) at every iteration.
The asymptotic notation \(\Theta(\cdot)\) in
\eqref{eq:spectral-scaling} is unrelated to the matrix normal
\(\Theta(W)\) defined below; the shared glyph should not be read as a shared
mathematical object.
Following Xie et al.~\cite{xie2026controlled}, we make the additional hard-constraint
modeling choice
\begin{equation}
 \mathcal L_R
 =
 \{W\in\R^{m\times n}:\snorm{W}=R\},
 \qquad
 R\asymp\sqrt{m/n}.
\tag{N2}
\label{eq:spectral-level-set}
\end{equation}
The normal \(\Theta\) is a consequence of this additional level-set model, not
an independent optimizer ingredient and not a consequence of \(\mu\)P alone.

\subsection{The smooth stratum and its unit normal}

Write the singular values of \(W\) as
\(\sigma_1\ge\sigma_2\ge\cdots\ge\sigma_q\), where
\(q=\min\{m,n\}\).  Throughout this subsection \(R>0\).
Statements below that invoke \(\sigma_2\), a second singular subspace, or a
two-leading-value eigensolver assume \(q\ge2\).  The one-dimensional case is
elementary and is excluded from those statements rather than hidden behind an
undefined \(\sigma_2\).

\begin{lemma}[Spectral-norm differential and generalized normal]
\label{lem:spectral-normal}
If \(\sigma_1(W)>\sigma_2(W)\), then the spectral norm is real analytic in a
neighborhood of \(W\), and
\begin{equation}
 D\snorm{W}[E]
 =
 \ip{u_1v_1^\top}{E}.
\tag{N3}
\label{eq:spectral-differential}
\end{equation}
Consequently
\begin{equation}
 \Theta(W):=u_1(W)v_1(W)^\top
\tag{N4}
\label{eq:true-normal}
\end{equation}
is the unique outward, gradient-oriented Frobenius unit normal to the smooth
stratum of \(\mathcal L_R\); the unoriented normal space is
\(\operatorname{span}\{\Theta(W)\}\).  It is invariant under the joint sign change
\((u_1,v_1)\mapsto(-u_1,-v_1)\), has rank one, and satisfies
\begin{equation}
 \snorm{\Theta}=\fnorm{\Theta}=\nnorm{\Theta}=1.
\tag{N5}
\end{equation}

If the positive leading singular value has multiplicity \(k>1\), then the
spectral norm is not differentiable and
\begin{equation}
 \partial\snorm{W}
 =
 \{U_kSV_k^\top:S\succeq0,\ \tr S=1\}.
\tag{N6}
\label{eq:set-valued-weight-normal}
\end{equation}
\end{lemma}

\begin{proof}
Apply simple-eigenvalue perturbation theory to the symmetric dilation
\[
 \mathcal H(W)
 =
 \begin{pmatrix}0&W\\W^\top&0\end{pmatrix}.
\]
Its largest eigenvalue is \(\sigma_1(W)\), with normalized eigenvector
\((u_1^\top,v_1^\top)^\top/\sqrt2\).  Differentiation gives
\eqref{eq:spectral-differential}.  Simplicity makes the one-dimensional left
and right singular spaces unique up to a joint sign, which cancels in the
outer product.  The norm identities follow because \(\Theta\) has one
nonzero singular value equal to one.  Formula
\eqref{eq:set-valued-weight-normal} is the standard spectral-norm
subdifferential formula \citep{watson1992subdifferential}.
\end{proof}

\begin{proposition}[Bouligand tangent cone at a repeated leading singular value]
\label{prop:repeated-top-tangent-cone}
Let \(W\in\mathcal L_R\) have a positive leading singular value \(R\) of
multiplicity \(k\), with corresponding matrices \(U_k,V_k\).  Then the
Bouligand tangent cone of the equality level set is
\begin{equation}
 T_{\mathcal L_R}^{\rm B}(W)
 =
 \left\{
 E:
 \lambda_{\max}\!\left(
 \operatorname{sym}(U_k^\top E V_k)
 \right)=0
 \right\},
\qquad
\operatorname{sym}(B)=\frac{B+B^\top}{2}.
\tag{N6a}
\label{eq:repeated-top-tangent-cone}
\end{equation}
For \(k=1\), this reduces to the hyperplane
\(\{E:\langle u_1v_1^\top,E\rangle_F=0\}\).  For \(k>1\), it is generally a
nonlinear, nonconvex cone and cannot be obtained by choosing one convenient
matrix from \(\partial\snorm{W}\).
\end{proposition}

\begin{proof}
The directional derivative of the largest singular value at a multiplicity
\(k\) cluster is
\begin{equation}
 \snorm{\cdot}'(W;E)
 =
 \lambda_{\max}\!\left(
 \operatorname{sym}(U_k^\top E V_k)
 \right).
\tag{N6b}
\label{eq:spectral-directional-derivative}
\end{equation}
This follows by compressing the perturbation of the symmetric dilation
\(\mathcal H(W)\) to its leading eigenspace.

For necessity, take any \(E\in T_{\mathcal L_R}^{\rm B}(W)\).  By definition
there are \(t_j\downarrow0\) and \(E_j\to E\) such that
\(W+t_jE_j\in\mathcal L_R\).  The spectral norm is Hadamard directionally
differentiable, hence
\[
 0
 =
 \lim_{j\to\infty}
 \frac{\snorm{W+t_jE_j}-\snorm{W}}{t_j}
 =
 \snorm{\cdot}'(W;E).
\]
For sufficiency, suppose the derivative in
\eqref{eq:spectral-directional-derivative} is zero and define
\[
 W(t)=R\,\frac{W+tE}{\snorm{W+tE}}.
\]
Then \(W(t)\in\mathcal L_R\) and
\(\snorm{W+tE}=R+o(t)\), so
\(W(t)=W+tE+o(t)\).  Thus \(E\) is a Bouligand tangent direction.
\end{proof}

\begin{boundary}
\textbf{The dual/weight symmetry gap.}
The manuscript treats rank deficiency of the \emph{dual matrix}
\(G+\lambda\Theta\) set-valuedly because that inclusion is the tractable inner
problem under study.  At a repeated leading singular value of the
\emph{weight}, it does not select one element of
\(\partial\snorm{W}\) and call the resulting hyperplane the tangent space.
\Cref{prop:repeated-top-tangent-cone} shows the actual first-order object.
The present inner-oracle analysis instead excludes this stratum from any
outer claim.  That is a genuine asymmetry of scope, not a resolution of the
global stratified problem; a tangent-cone/normal-cone outer method remains
open.
\end{boundary}

\begin{remark}[A normal is not a projector]
The matrix \(\Theta\) is a unit normal, not a projector.  Under the Frobenius metric the
orthogonal tangent projector is the linear map
\[
 \Pi_{T_W}(E)=E-\ip{\Theta}{E}\Theta.
\]
\end{remark}

\begin{lemma}[Dimension-free sensitivity of the normal]
\label{lem:normal-derivative}
Let \(\gamma_W=\sigma_1(W)-\sigma_2(W)>0\).  For every perturbation \(E\),
\begin{equation}
 \fnorm{D\Theta(W)[E]}
 \le
 \frac{\sqrt2}{\gamma_W}\snorm{E}.
\tag{N7}
\label{eq:normal-derivative}
\end{equation}
\end{lemma}

\begin{proof}
Choose differentiable singular-vector representatives satisfying
\(u_1^\top Du_1=v_1^\top Dv_1=0\).  The simple-eigenvector perturbation
 formula for \(\mathcal H(W)\) gives
\[
 \left\|D\left(\frac{u_1}{\sqrt2},\frac{v_1}{\sqrt2}\right)[E]\right\|_2
 \le \frac{\snorm{E}}{\gamma_W}.
\]
Since
\[
 D\Theta[E]=(Du_1[E])v_1^\top+u_1(Dv_1[E])^\top
\]
and the two summands are Frobenius orthogonal,
\[
 \fnorm{D\Theta[E]}^2
 =\|Du_1[E]\|_2^2+\|Dv_1[E]\|_2^2,
\]
which proves \eqref{eq:normal-derivative}.
\end{proof}

\subsection{Tangent constraint, ball oracle, and the role of retraction}

On the simple-leading-value stratum, the tangent space is the kernel of
\eqref{eq:spectral-differential}:
\begin{equation}
 T_W\mathcal L_R
 =
 \{\Phi:\ip{\Theta(W)}{\Phi}=0\}.
\tag{N8}
\label{eq:true-tangent-space}
\end{equation}
For a step \(E=-\eta R\Phi\), this also follows from
\begin{equation}
 \snorm{W-\eta R\Phi}
 =
 R-\eta R\ip{\Theta}{\Phi}
 +O\!\left(\frac{\eta^2R^2}{\gamma_W}\right).
\tag{N9}
\end{equation}
Thus the tangent LMO studied below is
\begin{equation}
 p^\star
 =
 \max_{\snorm{\Phi}\le1,\ \ip{\Theta}{\Phi}=0}
 \ip{G}{\Phi}.
\tag{P}
\label{eq:normal-primal}
\end{equation}
Here \(\Theta\) is computed from the current weight \(W\), whereas \(G\) is a
loss-gradient or momentum-like matrix supplied to the oracle.  They are
different inputs.  Their only coupling in the exact dual is through the affine
line \(G+\lambda\Theta\).

\begin{lemma}[Closed ball versus unit sphere]
\label{lem:ball-sphere}
If \(p^\star>0\), every optimizer of \eqref{eq:normal-primal} saturates
\(\snorm{\Phi}=1\); hence the ball and sphere formulations have the same
optimal value and the same optimizers.  If \(p^\star=0\), the optimal values
still agree provided \(mn\ge2\), but the ball formulation also contains
interior optima.  When \(mn=1\), the tangent hyperplane is \(\{0\}\), so the
unit-sphere problem has no feasible point.
\end{lemma}

\begin{proof}
If an optimizer with positive objective had spectral norm strictly below one,
positive rescaling would remain tangent and feasible while increasing the
objective.  If \(p^\star=0\), central symmetry of the feasible set implies
that \(\ip{G}{\Phi}=0\) for every tangent feasible \(\Phi\): otherwise either
\(\Phi\) or \(-\Phi\) would have positive objective.  When \(mn\ge2\), the
nonzero normal \(\Theta\) has a nonzero Frobenius-orthogonal complement.
Choose \(E\ne0\) in that complement and replace it by
\(E/\snorm{E}\).  This gives a tangent matrix of unit spectral norm, proving
equality of the optimal values.  The \(mn=1\) exception is immediate.
\end{proof}

This lemma explains why convex duality must be established on the closed ball,
even if an implementation later normalizes a nonzero accepted direction to
unit spectral norm.

\begin{lemma}[Explicit second-order drift]
\label{lem:normal-remainder}
Let \(\gamma_W>0\) and \(\snorm{E}\le\gamma_W/4\).  Then
\begin{equation}
 0
 \le
 \snorm{W+E}-\snorm{W}-\ip{\Theta}{E}
 \le
 \frac{2\snorm{E}^2}{\gamma_W}.
\tag{N10}
\label{eq:normal-remainder}
\end{equation}
In particular, a tangent step \(E=-\eta R\Phi\) with
\(\snorm{\Phi}\le1\) and \(\eta R\le\gamma_W/4\) has radial drift at most
\begin{equation}
 \frac{2\eta^2R^2}{\gamma_W}.
\tag{N11}
\label{eq:linearization-drift}
\end{equation}
\end{lemma}

\begin{proof}
Convexity gives the lower bound.  In a basis whose first vector is the leading
eigenvector of \(\mathcal H(W)\), write the symmetric dilation and its
perturbation as
\[
 \mathcal H(W)=\begin{pmatrix}\sigma_1&0\\0&B\end{pmatrix},
 \qquad
 \mathcal H(E)=\begin{pmatrix}a&b^\top\\b&C\end{pmatrix},
\]
with \(\lambda_{\max}(B)\le\sigma_1-\gamma_W\) and
\(\rho:=\snorm{E}=\snorm{\mathcal H(E)}\).  The Schur-complement equation for
the perturbed leading eigenvalue gives
\[
 0\le\widetilde\sigma_1-\sigma_1-a
 \le \frac{\rho^2}{\gamma_W-2\rho}
 \le \frac{2\rho^2}{\gamma_W}.
\]
Here \(a=\ip{\Theta}{E}\).  The dilation may have additional zero eigenvalues
when \(m\ne n\); their distance from \(\sigma_1\) is at least
\(\sigma_1\ge\gamma_W\), so the same bound applies.
\end{proof}

For \(W+E\ne0\), define the radial retraction
\begin{equation}
 \Retr_W(E)=R\frac{W+E}{\snorm{W+E}}.
\tag{N12}
\label{eq:radial-retraction}
\end{equation}
Its derivative is
\begin{equation}
 D\Retr_W(0)[E]
 =
 E-\frac WR\ip{\Theta}{E}.
\tag{N13}
\label{eq:retraction-derivative}
\end{equation}
Thus retraction guarantees exact radial feasibility, whereas the tangent LMO
chooses the best first-order direction under a prescribed spectral budget.
They are not interchangeable.  Retraction also does not automatically
preserve a fixed positive leading spectral gap.

\subsection{Approximate normals live in a projective geometry}

An implementation computes an approximate singular triplet
\((\hat\sigma,\hat u,\hat v)\), typically by a few power or block iterations,
and sets
\begin{equation}
 \hat\Theta=\hat u\hat v^\top.
\tag{N14}
\end{equation}
For classical bilateral power iteration with a starting vector having nonzero
leading-component overlap, the right-vector angle obeys
\begin{equation}
 \tan\angle(\hat v_k,v_1)
 \le
 \left(\frac{\sigma_2(W)}{\sigma_1(W)}\right)^{2k}
 \tan\angle(\hat v_0,v_1).
\tag{N14a}
\label{eq:power-rate}
\end{equation}
Thus a fixed iteration budget is reliable only when the leading ratio is
separated from one.

For a finite iteration budget, \eqref{eq:power-rate} applies with the actual
initial angle; a warm start changes that initial angle, not the exponent.
The equality \(\ip{\hat\Theta}{\Phi}=0\) is unchanged under
\(\hat\Theta\mapsto-\hat\Theta\).  The relevant object is therefore the normal
\emph{line}, not an oriented normal matrix.

Let
\begin{equation}
 \alpha_u=\arccos|u_1^\top\hat u|,
 \qquad
 \alpha_v=\arccos|v_1^\top\hat v|,
 \qquad
 \alpha_u,\alpha_v\in[0,\pi/2],
\tag{N15}
\label{eq:acute-angles}
\end{equation}
and put \(\vartheta=\alpha_u+\alpha_v\).

\begin{lemma}[Exact projective distance between rank-one normals]
\label{lem:projective-normal-distance}
\begin{equation}
 \min_{s\in\{\pm1\}}\nnorm{\Theta-s\hat\Theta}
 =
 2\sin\frac{\vartheta}{2},
\qquad
 \min_{s\in\{\pm1\}}\fnorm{\Theta-s\hat\Theta}
 =
 \sqrt{2-2\cos\alpha_u\cos\alpha_v}.
\tag{N16}
\label{eq:projective-distance}
\end{equation}
\end{lemma}

\begin{proof}
Choose \(s\) and a factorization
\(s\hat\Theta=\tilde u\tilde v^\top\) so that
\(u_1^\top\tilde u=\cos\alpha_u\ge0\) and
\(v_1^\top\tilde v=\cos\alpha_v\ge0\).  The difference has rank at most two.
The sum and product of its two squared singular values are
\[
 2-2\cos\alpha_u\cos\alpha_v,
 \qquad
 \sin^2\alpha_u\sin^2\alpha_v.
\]
Consequently the squared nuclear norm is
\[
 2-2\cos\alpha_u\cos\alpha_v
 +2\sin\alpha_u\sin\alpha_v
 =4\sin^2\frac{\alpha_u+\alpha_v}{2}.
\]
The other sign gives the larger of the two projective representatives.
\end{proof}

\begin{theorem}[Sharp worst-case leakage of an estimated tangent space]
\label{thm:sharp-normal-leakage}
For unit rank-one normals \(\Theta,\hat\Theta\),
\begin{equation}
 \sup_{\substack{\snorm{\Phi}\le1\\
                  \ip{\hat\Theta}{\Phi}=0}}
 |\ip{\Theta}{\Phi}|
 =
 \begin{cases}
 \sin\vartheta,&0\le\vartheta\le\pi/2,\\
 1,&\pi/2\le\vartheta\le\pi.
 \end{cases}
\tag{N17}
\label{eq:sharp-normal-leakage}
\end{equation}
\end{theorem}

\begin{proof}
The absolute value permits either sign of the objective.  After the projective
orientation used in the proof of \cref{lem:projective-normal-distance}, apply
T1 strong duality (Theorem~\ref{thm:T1}) to the oracle with objective \(\Theta\) and
normal \(\hat\Theta\):
\[
 \min_{\lambda\in\R}\nnorm{\Theta+\lambda\hat\Theta}.
\]
For \(\lambda=-t\), \(t\ge0\), the squared nuclear norm is
\[
 1+t^2-2t\cos(\alpha_u+\alpha_v).
\]
Its minimum is \(\sin^2\vartheta\) when
\(\cos\vartheta\ge0\), and \(1\) otherwise.  The branch
\(\lambda=t\ge0\) has squared norm
\[
 1+t^2+2t\cos(\alpha_u-\alpha_v),
\]
whose minimum is \(1\) because
\(|\alpha_u-\alpha_v|\le\pi/2\).  Comparing the branches proves
\eqref{eq:sharp-normal-leakage}.
\end{proof}

\begin{boundary}
\Cref{thm:sharp-normal-leakage} is a sharp \emph{uniform worst-case upper
bound}.  It is not a lower bound on every realized oracle direction: a
particular direction may be tangent to both normal hyperplanes.  Likewise, an
angle measured directly between two normal matrices is not, in general, the
sum \(\alpha_u+\alpha_v\) appearing in \eqref{eq:sharp-normal-leakage}.  No
empirical leakage percentage is inferred in this paper because the required
pair \((\alpha_u,\alpha_v)\) was not recorded.
\end{boundary}
